% 本文件由账本已入列事件生成；锚点首次呈现仅连接可核的年序与材料限度。
\section{南宋行在、边防与江南的加密年序}
以下补入按同年并列的政权与地区组织，而不是以一个朝廷的年号吞没其对手。每个可点击事件名保留账本 ID；叙述只将材料明确记录的动作置入年序，并把实施、规模、损失与动机留在可检验的限度内。

\subsection{1127年：南宋}
《宋史》为这一年留下了可回查的年序入口。

\noindent\textbf{南宋。}\eventanchor{SYD-1127-SS-001}{癸未，觀文殿大學士唐恪仰藥自殺}；\eventanchor{SYD-1127-SS-153}{何㮚、李若水勸帝親至軍中，從之，以太子監國而行}。

\subsection{1128年：南宋}
《宋史》为这一年留下了可回查的年序入口。

\noindent\textbf{南宋。}\eventanchor{SYD-1128-SS-002}{罷殿中侍御史馬伸，尋責濮州}；\eventanchor{SYD-1128-SS-154}{濱州賊蓋進陷棣州，守臣薑剛之死之}。

\subsection{1129年：南宋}
《宋史》为这一年留下了可回查的年序入口。

\noindent\textbf{南宋。}\eventanchor{SYD-1129-SS-003}{丙午，粘罕陷徐州，守臣王復及子倚死之，軍校趙立結鄉兵為興復計}；\eventanchor{SYD-1129-SS-155}{丙子，詔士民直言時政得失}。

\subsection{1130年：南宋}
《宋史》为这一年留下了可回查的年序入口。

\noindent\textbf{南宋。}\eventanchor{SYD-1130-SS-004}{罷諸路帥臣兼制置使、諸州守臣兼管內安撫使}；\eventanchor{SYD-1130-SS-156}{丙辰，復增左右司郎官為四員}。

\subsection{1131年：南宋、金与西夏（年序桥接）、南宋}
《宋史》、《金史》、《建炎以来系年要录》为这一年留下了可回查的年序入口。

\noindent\textbf{南宋。}\eventanchor{SYD-1131-SS-005}{丙辰，程昌寓遣杜湛擊楊華，敗之}；\eventanchor{SYD-1131-SS-157}{丙申，遣內侍撫問孔彥舟、桑仲}。

\noindent\textbf{南宋、金与西夏（年序桥接）。}\eventanchor{SY-1131-CHRONOLOGY-BRIDGE}{【C级年序桥接】保留1131 年的多政权并行检索入口；本条不主张所列政权在当年发生直接互动，具体事件须另以单项材料入账}。

\subsection{1132年：南宋}
《宋史》为这一年留下了可回查的年序入口。

\noindent\textbf{南宋。}\eventanchor{SYD-1132-SS-006}{班度量權衡于諸路，禁私造者}；\eventanchor{SYD-1132-SS-158}{丙辰，禁溫、台二州民結集社會}。

\subsection{1133年：南宋、金与西夏（年序桥接）、南宋}
《宋史》、《金史》、《建炎以来系年要录》为这一年留下了可回查的年序入口。

\noindent\textbf{南宋。}\eventanchor{SYD-1133-SS-007}{安化蠻犯邊，廣西經略使許中發兵擊之}；\eventanchor{SYD-1133-SS-159}{八月己丑，詔岳飛赴行在，留精兵萬人戍江州}。

\noindent\textbf{南宋、金与西夏（年序桥接）。}\eventanchor{SY-1133-CHRONOLOGY-BRIDGE}{【C级年序桥接】保留1133 年的多政权并行检索入口；本条不主张所列政权在当年发生直接互动，具体事件须另以单项材料入账}。

\subsection{1134年：南宋}
《宋史》为这一年留下了可回查的年序入口。

\noindent\textbf{南宋。}\eventanchor{SYD-1134-SS-008}{八月庚辰，以趙鼎知樞密院事，充川、陝宣撫處置使}；\eventanchor{SYD-1134-SS-160}{編次建炎以來詔旨，頒諸路}。

\subsection{1135年：南宋}
《宋史》为这一年留下了可回查的年序入口。

\noindent\textbf{南宋。}\eventanchor{SYD-1135-SS-009}{-\{岳\}-飛急攻湖賊水砦，賊將陳瑫降，楊太赴水死，餘黨劉衡等皆降}；\eventanchor{SYD-1135-SS-161}{-\{岳\}-飛為荊湖南北、襄陽府路制置使，將兵平湖賊楊太}。

\subsection{1136年：南宋、金与西夏（年序桥接）、南宋}
《宋史》、《金史》、《建炎以来系年要录》为这一年留下了可回查的年序入口。

\noindent\textbf{南宋。}\eventanchor{SYD-1136-SS-010}{-\{岳\}-飛及偽齊李成、孔彥舟連戰至蔡州，克之，偽守劉永壽舉城降}；\eventanchor{SYD-1136-SS-162}{-\{岳\}-飛遣統制牛皋破偽齊鎮汝軍，禽其守薛亨}。

\noindent\textbf{南宋、金与西夏（年序桥接）。}\eventanchor{SY-1136-CHRONOLOGY-BRIDGE}{【C级年序桥接】保留1136 年的多政权并行检索入口；本条不主张所列政权在当年发生直接互动，具体事件须另以单项材料入账}。

\subsection{1137年：南宋}
《宋史》为这一年留下了可回查的年序入口。

\noindent\textbf{南宋。}\eventanchor{SYD-1137-SS-011}{-\{岳\}-飛乞並統淮西兵以復京畿、陝右，許之，命飛盡護王德等諸將軍}；\eventanchor{SYD-1137-SS-163}{罷淮南提點司，東西兩路各置轉運兼提點刑獄、提舉茶鹽常平事}。

\subsection{1138年：南宋}
《宋史》为这一年留下了可回查的年序入口。

\noindent\textbf{南宋。}\eventanchor{SYD-1138-SS-012}{從官張燾、晏敦復、魏矼、曾開、李彌遜、尹焞、梁汝嘉、樓炤、蘇符、薛徽言、-\{御\}-史方廷實皆言不可}；\eventanchor{SYD-1138-SS-164}{八月戊午，詔：「日者遣使報聘鄰國，期還梓宮}。

\subsection{1139年：南宋}
《宋史》为这一年留下了可回查的年序入口。

\noindent\textbf{南宋。}\eventanchor{SYD-1139-SS-013}{八月己酉，復淮南諸州學官}；\eventanchor{SYD-1139-SS-165}{丙申，命詳驗劉豫偽官，換給告身}。

\subsection{1140年：南宋}
《宋史》为这一年留下了可回查的年序入口。

\noindent\textbf{南宋。}\eventanchor{SYD-1140-SS-014}{八月壬申朔，以張九成、喻樗、陳剛中、淩景夏、樊光遠、毛叔度、元盥等七人嘗不主和議，皆降黜之}；\eventanchor{SYD-1140-SS-166}{丙午，命兩浙、江東、福建諸州團結弓弩手}。

\subsection{1141年：南宋}
《宋史》为这一年留下了可回查的年序入口。

\noindent\textbf{南宋。}\eventanchor{SYD-1141-SS-015}{<u>李顯忠</u>遣統領<u>崔皋</u>擊敗<u>金</u>人于<u>舒城縣</u>}；\eventanchor{SYD-1141-SS-167}{<u>劉錡</u>自<u>東關</u>擊敗<u>金</u>人于<u>青谿</u>}。

\subsection{1142年：南宋}
《宋史》为这一年留下了可回查的年序入口。

\noindent\textbf{南宋。}\eventanchor{SYD-1142-SS-016}{八月辛酉朔，兀朮使來求商州及和尚、方山二原}；\eventanchor{SYD-1142-SS-168}{丙午，金使劉筈、完顏宗表等九人入見}。

\subsection{1143年：南宋、金与西夏（年序桥接）、南宋}
《宋史》、《金史》、《建炎以来系年要录》为这一年留下了可回查的年序入口。

\noindent\textbf{南宋。}\eventanchor{SYD-1143-SS-017}{八月丙戌，遣吏部侍郎江邈奉迎累朝神禦於溫州}；\eventanchor{SYD-1143-SS-169}{丙寅，處州兵士楊興等謀作亂，事覺伏誅}。

\noindent\textbf{南宋、金与西夏（年序桥接）。}\eventanchor{SY-1143-CHRONOLOGY-BRIDGE}{【C级年序桥接】保留1143 年的多政权并行检索入口；本条不主张所列政权在当年发生直接互动，具体事件须另以单项材料入账}。

\subsection{1144年：南宋}
《宋史》为这一年留下了可回查的年序入口。

\noindent\textbf{南宋。}\eventanchor{SYD-1144-SS-018}{丙午，罷-\{万俟\}-禼}；\eventanchor{SYD-1144-SS-170}{丙寅，立明法科兼經法}。

\subsection{1145年：南宋、金与西夏（年序桥接）、南宋}
《宋史》、《金史》、《建炎以来系年要录》为这一年留下了可回查的年序入口。

\noindent\textbf{南宋。}\eventanchor{SYD-1145-SS-019}{八月申戌朔，禁收折帛合零錢，止輸實數}；\eventanchor{SYD-1145-SS-171}{丙寅，全給秦檜歲賜公使錢萬緡}。

\noindent\textbf{南宋、金与西夏（年序桥接）。}\eventanchor{SY-1145-CHRONOLOGY-BRIDGE}{【C级年序桥接】保留1145 年的多政权并行检索入口；本条不主张所列政权在当年发生直接互动，具体事件须另以单项材料入账}。

\subsection{1146年：南宋、金与西夏（年序桥接）、南宋}
《宋史》、《金史》、《建炎以来系年要录》为这一年留下了可回查的年序入口。

\noindent\textbf{南宋。}\eventanchor{SYD-1146-SS-020}{丙申，復何鑄為端明殿學士兼侍讀}；\eventanchor{SYD-1146-SS-172}{丁亥，金遣烏古論海等來賀天申節}。

\noindent\textbf{南宋、金与西夏（年序桥接）。}\eventanchor{SY-1146-CHRONOLOGY-BRIDGE}{【C级年序桥接】保留1146 年的多政权并行检索入口；本条不主张所列政权在当年发生直接互动，具体事件须另以单项材料入账}。

\subsection{1147年：南宋、金与西夏（年序桥接）、南宋}
《宋史》、《金史》、《建炎以来系年要录》为这一年留下了可回查的年序入口。

\noindent\textbf{南宋。}\eventanchor{SYD-1147-SS-021}{八月庚子，罷建州創置賣鹽坊}；\eventanchor{SYD-1147-SS-173}{丙辰，金遣完顏宗藩等來賀明年正旦}。

\noindent\textbf{南宋、金与西夏（年序桥接）。}\eventanchor{SY-1147-CHRONOLOGY-BRIDGE}{【C级年序桥接】保留1147 年的多政权并行检索入口；本条不主张所列政权在当年发生直接互动，具体事件须另以单项材料入账}。

\subsection{1148年：南宋、金与西夏（年序桥接）、南宋}
《宋史》、《金史》、《建炎以来系年要录》为这一年留下了可回查的年序入口。

\noindent\textbf{南宋。}\eventanchor{SYD-1148-SS-022}{丙辰，加士夽開府儀同三司}；\eventanchor{SYD-1148-SS-174}{丙寅，借給被災農民春耕費}。

\noindent\textbf{南宋、金与西夏（年序桥接）。}\eventanchor{SY-1148-CHRONOLOGY-BRIDGE}{【C级年序桥接】保留1148 年的多政权并行检索入口；本条不主张所列政权在当年发生直接互动，具体事件须另以单项材料入账}。

\subsection{1149年：南宋}
《宋史》为这一年留下了可回查的年序入口。

\noindent\textbf{南宋。}\eventanchor{SYD-1149-SS-023}{丁丑，金遣完顏袞等來賀明年正旦}；\eventanchor{SYD-1149-SS-175}{丁未，減連、英、循、惠、新、恩六州免行錢}。

\subsection{1150年：南宋}
《宋史》为这一年留下了可回查的年序入口。

\noindent\textbf{南宋。}\eventanchor{SYD-1150-SS-024}{丙申，侍御史曹筠以附下罔上罷}；\eventanchor{SYD-1150-SS-176}{丙午，兩浙轉運副使曹泳言，李孟堅誦其父光所撰私史，語涉譏謗，詔送大理寺}。

\subsection{1151年：南宋、金与西夏（年序桥接）、南宋}
《宋史》、《金史》、《建炎以来系年要录》为这一年留下了可回查的年序入口。

\noindent\textbf{南宋。}\eventanchor{SYD-1151-SS-025}{八月辛未，秦檜上《重修諸路茶鹽法》}；\eventanchor{SYD-1151-SS-177}{丁亥，賜禮部進士趙逵以下四百四人及第、出身}。

\noindent\textbf{南宋、金与西夏（年序桥接）。}\eventanchor{SY-1151-CHRONOLOGY-BRIDGE}{【C级年序桥接】保留1151 年的多政权并行检索入口；本条不主张所列政权在当年发生直接互动，具体事件须另以单项材料入账}。

\subsection{1152年：南宋、金与西夏（年序桥接）、南宋}
《宋史》、《金史》、《建炎以来系年要录》为这一年留下了可回查的年序入口。

\noindent\textbf{南宋。}\eventanchor{SYD-1152-SS-026}{八月己卯，遣鄂州都統制田師中發兵同江西安撫使張澄、殿前司游奕軍統制李耕討述}；\eventanchor{SYD-1152-SS-178}{丁巳，立薦舉受財刑名}。

\noindent\textbf{南宋、金与西夏（年序桥接）。}\eventanchor{SY-1152-CHRONOLOGY-BRIDGE}{【C级年序桥接】保留1152 年的多政权并行检索入口；本条不主张所列政权在当年发生直接互动，具体事件须另以单项材料入账}。

\subsection{1153年：南宋}
《宋史》为这一年留下了可回查的年序入口。

\noindent\textbf{南宋。}\eventanchor{SYD-1153-SS-027}{八月乙丑，士撙薨，追封韶王}；\eventanchor{SYD-1153-SS-179}{丙寅，左宣教郎王孝廉謀據成都叛，事覺，伏誅}。

\subsection{1154年：南宋、金与西夏（年序桥接）、南宋}
《宋史》、《金史》、《建炎以来系年要录》为这一年留下了可回查的年序入口。

\noindent\textbf{南宋。}\eventanchor{SYD-1154-SS-028}{尋遣鐘世明如四川同議}；\eventanchor{SYD-1154-SS-180}{八月壬辰，禁百官避免輪對}。

\noindent\textbf{南宋、金与西夏（年序桥接）。}\eventanchor{SY-1154-CHRONOLOGY-BRIDGE}{【C级年序桥接】保留1154 年的多政权并行检索入口；本条不主张所列政权在当年发生直接互动，具体事件须另以单项材料入账}。

\subsection{1155年：南宋、金与西夏（年序桥接）、南宋}
《宋史》、《金史》、《建炎以来系年要录》为这一年留下了可回查的年序入口。

\noindent\textbf{南宋。}\eventanchor{SYD-1155-SS-029}{詔張浚、折彥質、-\{万俟\}-禼、段拂聽自便}；\eventanchor{SYD-1155-SS-181}{丙申，復以蕭振為四川制置使}。

\noindent\textbf{南宋、金与西夏（年序桥接）。}\eventanchor{SY-1155-CHRONOLOGY-BRIDGE}{【C级年序桥接】保留1155 年的多政权并行检索入口；本条不主张所列政权在当年发生直接互动，具体事件须另以单项材料入账}。

\subsection{1156年：南宋}
《宋史》为这一年留下了可回查的年序入口。

\noindent\textbf{南宋。}\eventanchor{SYD-1156-SS-030}{八月戊寅，班元豐、崇甯學制于諸路}；\eventanchor{SYD-1156-SS-182}{丙午，立互易薦舉坐罪法}。

\subsection{1157年：南宋、金与西夏（年序桥接）、南宋}
《宋史》、《金史》、《建炎以来系年要录》为这一年留下了可回查的年序入口。

\noindent\textbf{南宋。}\eventanchor{SYD-1157-SS-031}{丙辰，初命州縣置禁曆}；\eventanchor{SYD-1157-SS-183}{丙戌，賜禮部進士王十朋以下四百二十六人及第、出身}。

\noindent\textbf{南宋、金与西夏（年序桥接）。}\eventanchor{SY-1157-CHRONOLOGY-BRIDGE}{【C级年序桥接】保留1157 年的多政权并行检索入口；本条不主张所列政权在当年发生直接互动，具体事件须另以单项材料入账}。

\subsection{1158年：南宋、金与西夏（年序桥接）、南宋}
《宋史》、《金史》、《建炎以来系年要录》为这一年留下了可回查的年序入口。

\noindent\textbf{南宋。}\eventanchor{SYD-1158-SS-032}{己卯，命取公私銅器悉付鑄錢司，民間不輸者罪之}；\eventanchor{SYD-1158-SS-184}{八月戊子朔，置國史院，修神、哲、徽三朝正史}。

\noindent\textbf{南宋、金与西夏（年序桥接）。}\eventanchor{SY-1158-CHRONOLOGY-BRIDGE}{【C级年序桥接】保留1158 年的多政权并行检索入口；本条不主张所列政权在当年发生直接互动，具体事件须另以单项材料入账}。

\subsection{1159年：南宋、金与西夏（年序桥接）、南宋}
《宋史》、《金史》、《建炎以来系年要录》为这一年留下了可回查的年序入口。

\noindent\textbf{南宋。}\eventanchor{SYD-1159-SS-033}{八月甲子，募商人輸米行在諸倉，願以茶、鹽、礬鈔等償直者聽}；\eventanchor{SYD-1159-SS-185}{丙午，禁內外將佐營造、回易，掊斂軍士}。

\noindent\textbf{南宋、金与西夏（年序桥接）。}\eventanchor{SY-1159-CHRONOLOGY-BRIDGE}{【C级年序桥接】保留1159 年的多政权并行检索入口；本条不主张所列政权在当年发生直接互动，具体事件须另以单项材料入账}。

\subsection{1160年：南宋、金与西夏（年序桥接）、南宋}
《宋史》、《金史》、《建炎以来系年要录》为这一年留下了可回查的年序入口。

\noindent\textbf{南宋。}\eventanchor{SYD-1160-SS-034}{丙申，金遣蕭榮等來賀天申節}；\eventanchor{SYD-1160-SS-186}{丙申，命劉寶招制勝軍千人}。

\noindent\textbf{南宋、金与西夏（年序桥接）。}\eventanchor{SY-1160-CHRONOLOGY-BRIDGE}{【C级年序桥接】保留1160 年的多政权并行检索入口；本条不主张所列政权在当年发生直接互动，具体事件须另以单项材料入账}。

\subsection{1161年：南宋}
《宋史》为这一年留下了可回查的年序入口。

\noindent\textbf{南宋。}\eventanchor{SYD-1161-SS-035}{乙未，行新造會子於淮、浙、湖北、京西諸州}；\eventanchor{SYD-1161-SS-187}{八月辛丑朔，忠義人魏勝復海州，李寶承制以勝知州事}。

\subsection{1162年：南宋}
《宋史》为这一年留下了可回查的年序入口。

\noindent\textbf{南宋。}\eventanchor{SYD-1162-SS-036}{丙戌，給張浚錢十九萬緡，造沿江諸軍戰艦}；\eventanchor{SYD-1162-SS-188}{丙子，姚仲遣將復原州}。

\subsection{1163年：南宋}
《宋史》为这一年留下了可回查的年序入口。

\noindent\textbf{南宋。}\eventanchor{SYD-1163-SS-037}{安慶軍節度使士篯乞減奉賜之半，以助軍用}；\eventanchor{SYD-1163-SS-189}{八月丙寅，張浚復都督江、淮軍馬}。

\subsection{1164年：南宋}
《宋史》为这一年留下了可回查的年序入口。

\noindent\textbf{南宋。}\eventanchor{SYD-1164-SS-038}{罷兩浙、福建、江西、湖南、夔州路參議官}；\eventanchor{SYD-1164-SS-190}{丙申，命虞允文調兵討廣西諸盜}。

\subsection{1165年：南宋、金与西夏（年序桥接）、南宋}
《宋史》、《金史》、《建炎以来系年要录》为这一年留下了可回查的年序入口。

\noindent\textbf{南宋。}\eventanchor{SYD-1165-SS-039}{八月己卯，以永豐圩田賜建康都統司}；\eventanchor{SYD-1165-SS-191}{丙辰，詔有司治皇后家廟}。

\noindent\textbf{南宋、金与西夏（年序桥接）。}\eventanchor{SY-1165-CHRONOLOGY-BRIDGE}{【C级年序桥接】保留1165 年的多政权并行检索入口；本条不主张所列政权在当年发生直接互动，具体事件须另以单项材料入账}。

\subsection{1166年：南宋、金与西夏（年序桥接）、南宋}
《宋史》、《金史》、《建炎以来系年要录》为这一年留下了可回查的年序入口。

\noindent\textbf{南宋。}\eventanchor{SYD-1166-SS-040}{丁卯，賜禮部進士<u>蕭國梁</u>以下四百九十有三人及第、出身}；\eventanchor{SYD-1166-SS-192}{八月辛未朔，詔兩淮行鐵錢，銅錢毋過江北}。

\noindent\textbf{南宋、金与西夏（年序桥接）。}\eventanchor{SY-1166-CHRONOLOGY-BRIDGE}{【C级年序桥接】保留1166 年的多政权并行检索入口；本条不主张所列政权在当年发生直接互动，具体事件须另以单项材料入账}。

\subsection{1167年：南宋、金与西夏（年序桥接）、南宋}
《宋史》、《金史》、《建炎以来系年要录》为这一年留下了可回查的年序入口。

\noindent\textbf{南宋。}\eventanchor{SYD-1167-SS-041}{八月丁酉，內侍陳瑜、李宗回坐交結戚方受賂，瑜除名、決杖、黥面配循州，宗回除名、筠州編管，方責授果州團練副使、潭州安置，籍所盜庫金以犒軍}；\eventanchor{SYD-1167-SS-193}{罷成都、潼川路轉運司輪年銓試，以其事付制置司}。

\noindent\textbf{南宋、金与西夏（年序桥接）。}\eventanchor{SY-1167-CHRONOLOGY-BRIDGE}{【C级年序桥接】保留1167 年的多政权并行检索入口；本条不主张所列政权在当年发生直接互动，具体事件须另以单项材料入账}。

\subsection{1168年：南宋、金与西夏（年序桥接）、南宋}
《宋史》、《金史》、《建炎以来系年要录》为这一年留下了可回查的年序入口。

\noindent\textbf{南宋。}\eventanchor{SYD-1168-SS-042}{八月乙未，班祈雨雪之法于諸路}；\eventanchor{SYD-1168-SS-194}{丁亥，以饒、信二州、建甯府饑民嘯聚，遣官措置振濟}。

\noindent\textbf{南宋、金与西夏（年序桥接）。}\eventanchor{SY-1168-CHRONOLOGY-BRIDGE}{【C级年序桥接】保留1168 年的多政权并行检索入口；本条不主张所列政权在当年发生直接互动，具体事件须另以单项材料入账}。

\subsection{1169年：南宋、金与西夏（年序桥接）、南宋}
《宋史》、《金史》、《建炎以来系年要录》为这一年留下了可回查的年序入口。

\noindent\textbf{南宋。}\eventanchor{SYD-1169-SS-043}{罷利州路諸州營田官兵，募民耕佃}；\eventanchor{SYD-1169-SS-195}{丙寅，爲岳飛立廟于鄂州}。

\noindent\textbf{南宋、金与西夏（年序桥接）。}\eventanchor{SY-1169-CHRONOLOGY-BRIDGE}{【C级年序桥接】保留1169 年的多政权并行检索入口；本条不主张所列政权在当年发生直接互动，具体事件须另以单项材料入账}。

\subsection{1170年：南宋、金与西夏（年序桥接）、南宋}
《宋史》、《金史》、《建炎以来系年要录》为这一年留下了可回查的年序入口。

\noindent\textbf{南宋。}\eventanchor{SYD-1170-SS-044}{八月庚戌，虞允文請蚤建太子}；\eventanchor{SYD-1170-SS-196}{罷行在至鎮江徵稅所比近者十有三}。

\noindent\textbf{南宋、金与西夏（年序桥接）。}\eventanchor{SY-1170-CHRONOLOGY-BRIDGE}{【C级年序桥接】保留1170 年的多政权并行检索入口；本条不主张所列政权在当年发生直接互动，具体事件须另以单项材料入账}。

\subsection{1171年：南宋、金与西夏（年序桥接）、南宋}
《宋史》、《金史》、《建炎以来系年要录》为这一年留下了可回查的年序入口。

\noindent\textbf{南宋。}\eventanchor{SYD-1171-SS-045}{八月丙辰，詔兩淮民丁充民兵者，本名丁錢勿輸}；\eventanchor{SYD-1171-SS-197}{丙寅，金遣完顏宗甯等來賀明年正旦}。

\noindent\textbf{南宋、金与西夏（年序桥接）。}\eventanchor{SY-1171-CHRONOLOGY-BRIDGE}{【C级年序桥接】保留1171 年的多政权并行检索入口；本条不主张所列政权在当年发生直接互动，具体事件须另以单项材料入账}。

\subsection{1172年：南宋、金与西夏（年序桥接）、南宋}
《宋史》、《金史》、《建炎以来系年要录》为这一年留下了可回查的年序入口。

\noindent\textbf{南宋。}\eventanchor{SYD-1172-SS-046}{八年春正月庚午朔，班《乾道敕令格式》}；\eventanchor{SYD-1172-SS-198}{丙辰，金遣夾穀清臣等來賀會慶節}。

\noindent\textbf{南宋、金与西夏（年序桥接）。}\eventanchor{SY-1172-CHRONOLOGY-BRIDGE}{【C级年序桥接】保留1172 年的多政权并行检索入口；本条不主张所列政权在当年发生直接互动，具体事件须另以单项材料入账}。

\subsection{1173年：南宋}
《宋史》为这一年留下了可回查的年序入口。

\noindent\textbf{南宋。}\eventanchor{SYD-1173-SS-047}{八月丙子，詔興修水利}；\eventanchor{SYD-1173-SS-199}{丙辰，復分淮南安撫司爲東、西路}。

\subsection{1174年：南宋、金与西夏（年序桥接）、南宋}
《宋史》、《金史》、《建炎以来系年要录》为这一年留下了可回查的年序入口。

\noindent\textbf{南宋。}\eventanchor{SYD-1174-SS-048}{八月己未，張說罷爲太尉，提舉隆興府玉隆觀}；\eventanchor{SYD-1174-SS-200}{丙午，禁兩淮耕牛出境}。

\noindent\textbf{南宋、金与西夏（年序桥接）。}\eventanchor{SY-1174-CHRONOLOGY-BRIDGE}{【C级年序桥接】保留1174 年的多政权并行检索入口；本条不主张所列政权在当年发生直接互动，具体事件须另以单项材料入账}。

\subsection{1175年：南宋、金与西夏（年序桥接）、南宋}
《宋史》、《金史》、《建炎以来系年要录》为这一年留下了可回查的年序入口。

\noindent\textbf{南宋。}\eventanchor{SYD-1175-SS-049}{八月丙辰，江西總管賈和仲以捕茶寇失律，除名、賀州編管}；\eventanchor{SYD-1175-SS-201}{丁-\{丑\}-，遣左司諫湯邦彥等使金申議}。

\noindent\textbf{南宋、金与西夏（年序桥接）。}\eventanchor{SY-1175-CHRONOLOGY-BRIDGE}{【C级年序桥接】保留1175 年的多政权并行检索入口；本条不主张所列政权在当年发生直接互动，具体事件须另以单项材料入账}。

\subsection{1176年：南宋、金与西夏（年序桥接）、南宋}
《宋史》、《金史》、《建炎以来系年要录》为这一年留下了可回查的年序入口。

\noindent\textbf{南宋。}\eventanchor{SYD-1176-SS-050}{丁-\{丑\}-，命臨安守臣嚴禁逾侈}；\eventanchor{SYD-1176-SS-202}{丁酉，湯邦彥、陳雷奉使無狀，除名，邦彥新州、雷永州編管}。

\noindent\textbf{南宋、金与西夏（年序桥接）。}\eventanchor{SY-1176-CHRONOLOGY-BRIDGE}{【C级年序桥接】保留1176 年的多政权并行检索入口；本条不主张所列政权在当年发生直接互动，具体事件须另以单项材料入账}。

\subsection{1177年：南宋、金与西夏（年序桥接）、南宋}
《宋史》、《金史》、《建炎以来系年要录》为这一年留下了可回查的年序入口。

\noindent\textbf{南宋。}\eventanchor{SYD-1177-SS-051}{八月辛巳，禁耕牛過淮}；\eventanchor{SYD-1177-SS-203}{丁-\{丑\}-，詔監司、守臣歲舉武臣堪知縣者各二人}。

\noindent\textbf{南宋、金与西夏（年序桥接）。}\eventanchor{SY-1177-CHRONOLOGY-BRIDGE}{【C级年序桥接】保留1177 年的多政权并行检索入口；本条不主张所列政权在当年发生直接互动，具体事件须另以单项材料入账}。

\subsection{1178年：南宋、金与西夏（年序桥接）、南宋}
《宋史》、《金史》、《建炎以来系年要录》为这一年留下了可回查的年序入口。

\noindent\textbf{南宋。}\eventanchor{SYD-1178-SS-052}{八月甲午，詔諸路監司戒所部，民稅毋以重價強折輸錢}；\eventanchor{SYD-1178-SS-204}{丙辰，金遣烏延察等來賀明年正旦}。

\noindent\textbf{南宋、金与西夏（年序桥接）。}\eventanchor{SY-1178-CHRONOLOGY-BRIDGE}{【C级年序桥接】保留1178 年的多政权并行检索入口；本条不主张所列政权在当年发生直接互动，具体事件须另以单项材料入账}。

\subsection{1179年：南宋、金与西夏（年序桥接）、南宋}
《宋史》、《金史》、《建炎以来系年要录》为这一年留下了可回查的年序入口。

\noindent\textbf{南宋。}\eventanchor{SYD-1179-SS-053}{八月庚寅，罷諸路監司、帥守便宜行事}；\eventanchor{SYD-1179-SS-205}{丙申，詔太學兩優釋褐，與殿試第二人恩例}。

\noindent\textbf{南宋、金与西夏（年序桥接）。}\eventanchor{SY-1179-CHRONOLOGY-BRIDGE}{【C级年序桥接】保留1179 年的多政权并行检索入口；本条不主张所列政权在当年发生直接互动，具体事件须另以单项材料入账}。

\subsection{1180年：南宋、金与西夏（年序桥接）、南宋}
《宋史》、《金史》、《建炎以来系年要录》为这一年留下了可回查的年序入口。

\noindent\textbf{南宋。}\eventanchor{SYD-1180-SS-054}{戊子，遣葉宏等使金賀正旦}；\eventanchor{SYD-1180-SS-206}{八月癸未，禁黎州官吏市蕃商物}。

\noindent\textbf{南宋、金与西夏（年序桥接）。}\eventanchor{SY-1180-CHRONOLOGY-BRIDGE}{【C级年序桥接】保留1180 年的多政权并行检索入口；本条不主张所列政权在当年发生直接互动，具体事件须另以单项材料入账}。

\subsection{1181年：南宋、金与西夏（年序桥接）、南宋}
《宋史》、《金史》、《建炎以来系年要录》为这一年留下了可回查的年序入口。

\noindent\textbf{南宋。}\eventanchor{SYD-1181-SS-055}{八月丙午，以旱，罷招軍}；\eventanchor{SYD-1181-SS-207}{丙辰，以臨安疫，分命醫官診視軍民}。

\noindent\textbf{南宋、金与西夏（年序桥接）。}\eventanchor{SY-1181-CHRONOLOGY-BRIDGE}{【C级年序桥接】保留1181 年的多政权并行检索入口；本条不主张所列政权在当年发生直接互动，具体事件须另以单项材料入账}。

\subsection{1182年：南宋、金与西夏（年序桥接）、南宋}
《宋史》、《金史》、《建炎以来系年要录》为这一年留下了可回查的年序入口。

\noindent\textbf{南宋。}\eventanchor{SYD-1182-SS-056}{八月己亥朔，詔紹興民戶去歲已納夏稅應減者三十萬緡，理爲今年之數}；\eventanchor{SYD-1182-SS-208}{丙子，以子彤爲容州觀察使，封安定郡王}。

\noindent\textbf{南宋、金与西夏（年序桥接）。}\eventanchor{SY-1182-CHRONOLOGY-BRIDGE}{【C级年序桥接】保留1182 年的多政权并行检索入口；本条不主张所列政权在当年发生直接互动，具体事件须另以单项材料入账}。

\subsection{1183年：南宋、金与西夏（年序桥接）、南宋}
《宋史》、《金史》、《建炎以来系年要录》为这一年留下了可回查的年序入口。

\noindent\textbf{南宋。}\eventanchor{SYD-1183-SS-057}{丁-\{丑\}-，詔除災傷州縣淳熙八年欠稅}；\eventanchor{SYD-1183-SS-209}{丁亥，金遣完顏婆盧火等來賀明年正旦}。

\noindent\textbf{南宋、金与西夏（年序桥接）。}\eventanchor{SY-1183-CHRONOLOGY-BRIDGE}{【C级年序桥接】保留1183 年的多政权并行检索入口；本条不主张所列政权在当年发生直接互动，具体事件须另以单项材料入账}。

\subsection{1184年：南宋、金与西夏（年序桥接）、南宋}
《宋史》、《金史》、《建炎以来系年要录》为这一年留下了可回查的年序入口。

\noindent\textbf{南宋。}\eventanchor{SYD-1184-SS-058}{八月庚申，遣章森使金賀正旦}；\eventanchor{SYD-1184-SS-210}{丙午，詔江東、西路諸監司，義役、差役從民便}。

\noindent\textbf{南宋、金与西夏（年序桥接）。}\eventanchor{SY-1184-CHRONOLOGY-BRIDGE}{【C级年序桥接】保留1184 年的多政权并行检索入口；本条不主张所列政权在当年发生直接互动，具体事件须另以单项材料入账}。

\subsection{1185年：南宋、金与西夏（年序桥接）、南宋}
《宋史》、《金史》、《建炎以来系年要录》为这一年留下了可回查的年序入口。

\noindent\textbf{南宋。}\eventanchor{SYD-1185-SS-059}{八月癸亥，詔太上皇壽八十，令有司議慶壽禮}；\eventanchor{SYD-1185-SS-211}{丙戌，詔恤潮州、台州被水之家}。

\noindent\textbf{南宋、金与西夏（年序桥接）。}\eventanchor{SY-1185-CHRONOLOGY-BRIDGE}{【C级年序桥接】保留1185 年的多政权并行检索入口；本条不主张所列政权在当年发生直接互动，具体事件须另以单项材料入账}。

\subsection{1186年：南宋、金与西夏（年序桥接）、南宋}
《宋史》、《金史》、《建炎以来系年要录》为这一年留下了可回查的年序入口。

\noindent\textbf{南宋。}\eventanchor{SYD-1186-SS-060}{丙申，賜沖晦處士郭雍號曰頤正先生，仍遣官就問雍所欲言，備錄來上}；\eventanchor{SYD-1186-SS-212}{丙寅，梁克家罷爲觀文殿大學士、醴泉觀使兼侍讀}。

\noindent\textbf{南宋、金与西夏（年序桥接）。}\eventanchor{SY-1186-CHRONOLOGY-BRIDGE}{【C级年序桥接】保留1186 年的多政权并行检索入口；本条不主张所列政权在当年发生直接互动，具体事件须另以单项材料入账}。

\subsection{1187年：南宋、金与西夏（年序桥接）、南宋}
《宋史》、《金史》、《建炎以来系年要录》为这一年留下了可回查的年序入口。

\noindent\textbf{南宋。}\eventanchor{SYD-1187-SS-061}{八月辛未，賜度牒一百道、米四萬五千石，備振紹興府饑}；\eventanchor{SYD-1187-SS-213}{丙辰，命臨安府捕蝗，募民輸米振濟}。

\noindent\textbf{南宋、金与西夏（年序桥接）。}\eventanchor{SY-1187-CHRONOLOGY-BRIDGE}{【C级年序桥接】保留1187 年的多政权并行检索入口；本条不主张所列政权在当年发生直接互动，具体事件须另以单项材料入账}。

\subsection{1188年：南宋、金与西夏（年序桥接）、南宋}
《宋史》、《金史》、《建炎以来系年要录》为这一年留下了可回查的年序入口。

\noindent\textbf{南宋。}\eventanchor{SYD-1188-SS-062}{丁巳，詔修《高宗實錄》}；\eventanchor{SYD-1188-SS-214}{己-\{丑\}-，詔減臨安、紹興府囚罪一等，釋杖以下，民緣欑宮役者蠲其賦}。

\noindent\textbf{南宋、金与西夏（年序桥接）。}\eventanchor{SY-1188-CHRONOLOGY-BRIDGE}{【C级年序桥接】保留1188 年的多政权并行检索入口；本条不主张所列政权在当年发生直接互动，具体事件须另以单项材料入账}。

\subsection{1189年：南宋}
《宋史》为这一年留下了可回查的年序入口。

\noindent\textbf{南宋。}\eventanchor{SYD-1189-SS-063}{甲午，封孫抦爲嘉國公}；\eventanchor{SYD-1189-SS-215}{壬戌，下詔傳位皇太子}。

\subsection{1190年：南宋、金与西夏（年序桥接）、南宋}
《宋史》、《金史》、《建炎以来系年要录》为这一年留下了可回查的年序入口。

\noindent\textbf{南宋。}\eventanchor{SYD-1190-SS-064}{八月辛卯，立任子中銓人吏部簾試法}；\eventanchor{SYD-1190-SS-216}{丙申，以上供等錢償廣州放免身丁錢數}。

\noindent\textbf{南宋、金与西夏（年序桥接）。}\eventanchor{SY-1190-CHRONOLOGY-BRIDGE}{【C级年序桥接】保留1190 年的多政权并行检索入口；本条不主张所列政权在当年发生直接互动，具体事件须另以单项材料入账}。

\subsection{1191年：南宋、金与西夏（年序桥接）、南宋}
《宋史》、《金史》、《建炎以来系年要录》为这一年留下了可回查的年序入口。

\noindent\textbf{南宋。}\eventanchor{SYD-1191-SS-065}{光宗紹熙二年，加諡受命中興全功至德聖神武文昭仁憲孝皇帝}；\eventanchor{SYD-1191-SS-217}{丙申，詔侍從、兩省、台諫及在外侍從之臣，各舉所知嘗任監司、郡守可充郎官、卿監及資歷未深可充諸職事官者各三人}。

\noindent\textbf{南宋、金与西夏（年序桥接）。}\eventanchor{SY-1191-CHRONOLOGY-BRIDGE}{【C级年序桥接】保留1191 年的多政权并行检索入口；本条不主张所列政权在当年发生直接互动，具体事件须另以单项材料入账}。

\subsection{1192年：南宋、金与西夏（年序桥接）、南宋}
《宋史》、《金史》、《建炎以来系年要录》为这一年留下了可回查的年序入口。

\noindent\textbf{南宋。}\eventanchor{SYD-1192-SS-066}{八月甲寅，詔兩淮行鐵錢交子}；\eventanchor{SYD-1192-SS-218}{兵部尚書羅點、給事中尤袤、中書舍人黃裳皆上疏請帝朝重華宮，吏部尚書趙汝愚亦因面對以請，帝開納}。

\noindent\textbf{南宋、金与西夏（年序桥接）。}\eventanchor{SY-1192-CHRONOLOGY-BRIDGE}{【C级年序桥接】保留1192 年的多政权并行检索入口；本条不主张所列政权在当年发生直接互动，具体事件须另以单项材料入账}。

\subsection{1193年：南宋}
《宋史》为这一年留下了可回查的年序入口。

\noindent\textbf{南宋。}\eventanchor{SYD-1193-SS-067}{明日會慶節，帝以疾不果朝，丞相葛邲率百官賀于重華宮}；\eventanchor{SYD-1193-SS-219}{丙寅，貸淮西民市牛錢}。

\subsection{1194年：南宋}
《宋史》为这一年留下了可回查的年序入口。

\noindent\textbf{南宋。}\eventanchor{SYD-1194-SS-068}{丙辰，侍講黃裳、秘書少監孫逢吉等再上疏以請}；\eventanchor{SYD-1194-SS-220}{丙戌，權戶部侍郎袁說友入對，請朝重華宮}。

\subsection{1195年：南宋、金与西夏（年序桥接）、南宋}
《宋史》、《金史》、《建炎以来系年要录》为这一年留下了可回查的年序入口。

\noindent\textbf{南宋。}\eventanchor{SYD-1195-SS-069}{八月己巳，詔內外諸軍主帥條奏武備邊防之策以聞}；\eventanchor{SYD-1195-SS-221}{丙午，以監察-\{御\}-史胡紘言，責授趙汝愚寧遠軍節度副使、永州安置}。

\noindent\textbf{南宋、金与西夏（年序桥接）。}\eventanchor{SY-1195-CHRONOLOGY-BRIDGE}{【C级年序桥接】保留1195 年的多政权并行检索入口；本条不主张所列政权在当年发生直接互动，具体事件须另以单项材料入账}。

\subsection{1196年：南宋}
《宋史》为这一年留下了可回查的年序入口。

\noindent\textbf{南宋。}\eventanchor{SYD-1196-SS-070}{丙辰，乙太常少卿胡紘請，權住進擬偽學之黨}；\eventanchor{SYD-1196-SS-222}{丙戌，減諸路死罪囚，釋流以下}。

\subsection{1197年：南宋、金与西夏（年序桥接）、南宋}
《宋史》、《金史》、《建炎以来系年要录》为这一年留下了可回查的年序入口。

\noindent\textbf{南宋。}\eventanchor{SYD-1197-SS-071}{八月戊子，復置嚴州神泉監}；\eventanchor{SYD-1197-SS-223}{丙戌，金遣完顏愈來賀瑞慶節}。

\noindent\textbf{南宋、金与西夏（年序桥接）。}\eventanchor{SY-1197-CHRONOLOGY-BRIDGE}{【C级年序桥接】保留1197 年的多政权并行检索入口；本条不主张所列政权在当年发生直接互动，具体事件须另以单项材料入账}。

\subsection{1198年：南宋、金与西夏（年序桥接）、南宋}
《宋史》、《金史》、《建炎以来系年要录》为这一年留下了可回查的年序入口。

\noindent\textbf{南宋。}\eventanchor{SYD-1198-SS-072}{八月丁卯朔，以久雨，決系囚}；\eventanchor{SYD-1198-SS-224}{丙戌，詔以太上皇聖躬清復，率群臣上壽}。

\noindent\textbf{南宋、金与西夏（年序桥接）。}\eventanchor{SY-1198-CHRONOLOGY-BRIDGE}{【C级年序桥接】保留1198 年的多政权并行检索入口；本条不主张所列政权在当年发生直接互动，具体事件须另以单项材料入账}。

\subsection{1199年：南宋、金与西夏（年序桥接）、南宋}
《宋史》、《金史》、《建炎以来系年要录》为这一年留下了可回查的年序入口。

\noindent\textbf{南宋。}\eventanchor{SYD-1199-SS-073}{丙辰，遣朱致知使金賀正旦}；\eventanchor{SYD-1199-SS-225}{丙戌，詔減諸路流囚，釋杖以下，推恩如慶壽故事}。

\noindent\textbf{南宋、金与西夏（年序桥接）。}\eventanchor{SY-1199-CHRONOLOGY-BRIDGE}{【C级年序桥接】保留1199 年的多政权并行检索入口；本条不主张所列政权在当年发生直接互动，具体事件须另以单项材料入账}。

\subsection{1200年：南宋、金与西夏（年序桥接）、南宋}
《宋史》、《金史》、《建炎以来系年要录》为这一年留下了可回查的年序入口。

\noindent\textbf{南宋。}\eventanchor{SYD-1200-SS-074}{丙寅，詔大理、三衙、臨安府及諸路闕雨州縣釋杖以下囚}；\eventanchor{SYD-1200-SS-226}{丁-\{丑\}-，詔三省、樞密院擇臣僚封事可行者以聞}。

\noindent\textbf{南宋、金与西夏（年序桥接）。}\eventanchor{SY-1200-CHRONOLOGY-BRIDGE}{【C级年序桥接】保留1200 年的多政权并行检索入口；本条不主张所列政权在当年发生直接互动，具体事件须另以单项材料入账}。

\subsection{1201年：南宋、金与西夏（年序桥接）、南宋}
《宋史》、《金史》、《建炎以来系年要录》为这一年留下了可回查的年序入口。

\noindent\textbf{南宋。}\eventanchor{SYD-1201-SS-075}{頒《慶元寬恤詔令》、《役法撮要》}；\eventanchor{SYD-1201-SS-227}{丙午，詔文武臣無寓居州任厘務官，著爲令}。

\noindent\textbf{南宋、金与西夏（年序桥接）。}\eventanchor{SY-1201-CHRONOLOGY-BRIDGE}{【C级年序桥接】保留1201 年的多政权并行检索入口；本条不主张所列政权在当年发生直接互动，具体事件须另以单项材料入账}。

\subsection{1202年：南宋、金与西夏（年序桥接）、南宋}
《宋史》、《金史》、《建炎以来系年要录》为这一年留下了可回查的年序入口。

\noindent\textbf{南宋。}\eventanchor{SYD-1202-SS-076}{丙寅，嗣秀王伯圭薨，追封崇王，諡曰憲靖}；\eventanchor{SYD-1202-SS-228}{丁亥，修《高宗正史》、《寶訓》}。

\noindent\textbf{南宋、金与西夏（年序桥接）。}\eventanchor{SY-1202-CHRONOLOGY-BRIDGE}{【C级年序桥接】保留1202 年的多政权并行检索入口；本条不主张所列政权在当年发生直接互动，具体事件须另以单项材料入账}。

\subsection{1203年：南宋、金与西夏（年序桥接）、南宋}
《宋史》、《金史》、《建炎以来系年要录》为这一年留下了可回查的年序入口。

\noindent\textbf{南宋。}\eventanchor{SYD-1203-SS-077}{八月壬寅，增置襄陽騎軍}；\eventanchor{SYD-1203-SS-229}{丙午，出封樁庫兩淮交子一百萬，命轉運司收民間鐵錢}。

\noindent\textbf{南宋、金与西夏（年序桥接）。}\eventanchor{SY-1203-CHRONOLOGY-BRIDGE}{【C级年序桥接】保留1203 年的多政权并行检索入口；本条不主张所列政权在当年发生直接互动，具体事件须另以单项材料入账}。

\subsection{1204年：南宋}
《宋史》为这一年留下了可回查的年序入口。

\noindent\textbf{南宋。}\eventanchor{SYD-1204-SS-078}{丙辰，除靜江府、昭州折布錢}；\eventanchor{SYD-1204-SS-230}{丙申，置諸軍帳前雄校，以軍官子孫補之}。

\subsection{1205年：南宋、金与西夏（年序桥接）、南宋}
《宋史》、《金史》、《建炎以来系年要录》为这一年留下了可回查的年序入口。

\noindent\textbf{南宋。}\eventanchor{SYD-1205-SS-079}{八月丙戌朔，蠲兩浙闕雨州縣贓賞錢}；\eventanchor{SYD-1205-SS-231}{丙寅，升嘉定府爲嘉慶軍}。

\noindent\textbf{南宋、金与西夏（年序桥接）。}\eventanchor{SY-1205-CHRONOLOGY-BRIDGE}{【C级年序桥接】保留1205 年的多政权并行检索入口；本条不主张所列政权在当年发生直接互动，具体事件须另以单项材料入账}。

\subsection{1206年：南宋}
《宋史》为这一年留下了可回查的年序入口。

\noindent\textbf{南宋。}\eventanchor{SYD-1206-SS-080}{八月丙寅，有司上《開禧刑名斷例》}；\eventanchor{SYD-1206-SS-232}{罷旱傷州軍比較租賦一年}。

\subsection{1207年：南宋}
《宋史》为这一年留下了可回查的年序入口。

\noindent\textbf{南宋。}\eventanchor{SYD-1207-SS-081}{辛巳，再奪鄧友龍五官、南雄州安置，尋除名，徙循州}；\eventanchor{SYD-1207-SS-233}{丙戌，命淮西轉運司措置雄淮軍}。

\subsection{1208年：南宋}
《宋史》为这一年留下了可回查的年序入口。

\noindent\textbf{南宋。}\eventanchor{SYD-1208-SS-082}{八月戊辰朔，發米振貧民}；\eventanchor{SYD-1208-SS-234}{丙申，幸太乙宮、明慶寺禱雨}。

\subsection{1209年：南宋}
《宋史》为这一年留下了可回查的年序入口。

\noindent\textbf{南宋。}\eventanchor{SYD-1209-SS-083}{丙戌，發米十萬石振兩淮饑民}；\eventanchor{SYD-1209-SS-235}{丙戌，金遣使來賀明年正旦}。

\subsection{1210年：南宋}
《宋史》为这一年留下了可回查的年序入口。

\noindent\textbf{南宋。}\eventanchor{SYD-1210-SS-084}{丙辰，以久雨，釋兩浙州縣繫囚}；\eventanchor{SYD-1210-SS-236}{丙寅，湖南賊羅世傳縛李元礪以降，峒寇悉平}。

\subsection{1211年：南宋}
《宋史》为这一年留下了可回查的年序入口。

\noindent\textbf{南宋。}\eventanchor{SYD-1211-SS-085}{丙午，賜黑風峒名曰效忠}；\eventanchor{SYD-1211-SS-237}{丙午，詔湖南、江西諸州經賊蹂踐者，監司、守臣考縣令安集之實，第其能否以聞}。

\subsection{1212年：南宋}
《宋史》为这一年留下了可回查的年序入口。

\noindent\textbf{南宋。}\eventanchor{SYD-1212-SS-086}{八月甲戌朔，-\{御\}-後殿，復膳}；\eventanchor{SYD-1212-SS-238}{冬十月辛巳，詔諸路總領官歲舉堪將帥者二三人，安撫、提刑舉可備將材者各二人}。

\subsection{1213年：南宋}
《宋史》为这一年留下了可回查的年序入口。

\noindent\textbf{南宋。}\eventanchor{SYD-1213-SS-087}{八月己巳朔，詔諸路監司、帥臣舉所部官吏之才行卓絕、績用章著者}；\eventanchor{SYD-1213-SS-239}{丙申，以雷發非時，下罪己詔}。

\subsection{1214年：南宋}
《宋史》为这一年留下了可回查的年序入口。

\noindent\textbf{南宋。}\eventanchor{SYD-1214-SS-088}{八月癸巳朔，罷關外四州所增方田稅}；\eventanchor{SYD-1214-SS-240}{罷四川制置大使司所開鹽井}。

\subsection{1215年：南宋}
《宋史》为这一年留下了可回查的年序入口。

\noindent\textbf{南宋。}\eventanchor{SYD-1215-SS-089}{八年春正月辛未，命師禹嗣秀王}；\eventanchor{SYD-1215-SS-241}{八月己-\{丑\}-，賜張栻諡曰宣}。

\subsection{1216年：南宋}
《宋史》为这一年留下了可回查的年序入口。

\noindent\textbf{南宋。}\eventanchor{SYD-1216-SS-090}{丙寅，金遣使來賀瑞慶節}；\eventanchor{SYD-1216-SS-242}{丙子，命諸州招填軍籍}。

\subsection{1217年：南宋}
《宋史》为这一年留下了可回查的年序入口。

\noindent\textbf{南宋。}\eventanchor{SYD-1217-SS-091}{八月乙丑，詔監司、郡守各舉威勇才略可將帥者二人}；\eventanchor{SYD-1217-SS-243}{丙辰，詔江淮制置使李玨、京湖制置使趙方措置調遣，仍聽便宜行事}。

\subsection{1218年：南宋}
《宋史》为这一年留下了可回查的年序入口。

\noindent\textbf{南宋。}\eventanchor{SYD-1218-SS-092}{丁亥，詔侍從、臺諫、兩省官集議平戎、禦戎、和戎三策}；\eventanchor{SYD-1218-SS-244}{丁酉，詔四川忠義人立功，賞視官軍}。

\subsection{1219年：南宋}
《宋史》为这一年留下了可回查的年序入口。

\noindent\textbf{南宋。}\eventanchor{SYD-1219-SS-093}{八月戊辰，復合利州東、西路爲一}；\eventanchor{SYD-1219-SS-245}{丁-\{丑\}-，雅州蠻入盧山縣}。

\subsection{1220年：南宋}
《宋史》为这一年留下了可回查的年序入口。

\noindent\textbf{南宋。}\eventanchor{SYD-1220-SS-094}{丙辰，四川宣撫司招黎人土丁，降之}；\eventanchor{SYD-1220-SS-246}{丁巳，黎州土丁叛，遣兵討之}。

\subsection{1221年：南宋}
《宋史》为这一年留下了可回查的年序入口。

\noindent\textbf{南宋。}\eventanchor{SYD-1221-SS-095}{八月乙卯，賜史彌遠家廟}；\eventanchor{SYD-1221-SS-247}{丙寅，夏人復以書來四川趣會兵}。

\subsection{1222年：南宋}
《宋史》为这一年留下了可回查的年序入口。

\noindent\textbf{南宋。}\eventanchor{SYD-1222-SS-096}{八月己卯，命戶部詳議義役}；\eventanchor{SYD-1222-SS-248}{丁巳，復以隨州三關隸德安府，置關使}。

\subsection{1223年：南宋}
《宋史》为这一年留下了可回查的年序入口。

\noindent\textbf{南宋。}\eventanchor{SYD-1223-SS-097}{丁卯，以道州民饑，詔發米振之}；\eventanchor{SYD-1223-SS-249}{冬十一月辛亥，以太平州大水，詔振恤之}。

\subsection{1224年：南宋}
《宋史》为这一年留下了可回查的年序入口。

\noindent\textbf{南宋。}\eventanchor{SYD-1224-SS-098}{八月乙亥，罷通州天賜鹽場}；\eventanchor{SYD-1224-SS-250}{癸亥，命淮東西、湖北路轉運司提督營屯田}。

\subsection{1225年：南宋}
《宋史》为这一年留下了可回查的年序入口。

\noindent\textbf{南宋。}\eventanchor{SYD-1225-SS-099}{丙寅，詔不熄爲保康軍承宣使、嗣濮王}；\eventanchor{SYD-1225-SS-251}{壬寅，帝兩請皇太后垂簾，不允}。

\subsection{1226年：南宋}
《宋史》为这一年留下了可回查的年序入口。

\noindent\textbf{南宋。}\eventanchor{SYD-1226-SS-100}{八月乙巳，濟王竑追降巴陵縣公}；\eventanchor{SYD-1226-SS-252}{冬十月甲申，詔《寧宗-\{御\}-集》閣以「寶章」爲名，仍置學士、待制員}。

\subsection{1227年：南宋}
《宋史》为这一年留下了可回查的年序入口。

\noindent\textbf{南宋。}\eventanchor{SYD-1227-SS-101}{己巳，詔：「朕觀朱熹集注《大學》、《論語》、《孟子》、《中庸》，發揮聖賢蘊奧，有補治道，朕勵志講學，緬懷典刑，可特贈熹太師，追封信國公}；\eventanchor{SYD-1227-SS-253}{八月庚戌，詔謝氏特封通義郡夫人}。

\subsection{1228年：南宋}
《宋史》为这一年留下了可回查的年序入口。

\noindent\textbf{南宋。}\eventanchor{SYD-1228-SS-102}{丁酉，詔申嚴皇城司給符之制，照闌入法}。

\subsection{1229年：南宋}
《宋史》为这一年留下了可回查的年序入口。

\noindent\textbf{南宋。}\eventanchor{SYD-1229-SS-103}{冬十月壬戌，詔台州水災，除民田租及茶、鹽、酒酤諸雜稅，郡縣抑納者監司察之}；\eventanchor{SYD-1229-SS-254}{二月庚戌，詔歲舉廉吏或犯奸贓，保任同坐，監司、守臣其申嚴覺察}。

\subsection{1230年：南宋}
《宋史》为这一年留下了可回查的年序入口。

\noindent\textbf{南宋。}\eventanchor{SYD-1230-SS-104}{丁卯，冊命貴妃謝氏爲皇后}；\eventanchor{SYD-1230-SS-255}{乙丑，詔免明年元會禮}。

\subsection{1231年：南宋}
《宋史》为这一年留下了可回查的年序入口。

\noindent\textbf{南宋。}\eventanchor{SYD-1231-SS-105}{八月己未，大元兵破武休，入興元，攻仙人關}；\eventanchor{SYD-1231-SS-256}{丙子，詔起復孟珙從義郎、京西路分，棗陽軍駐紮}。

\subsection{1232年：南宋}
《宋史》为这一年留下了可回查的年序入口。

\noindent\textbf{南宋。}\eventanchor{SYD-1232-SS-106}{壬辰，史嵩之進大理卿、權刑部侍郎、京湖安撫制置使、知襄陽府}；\eventanchor{SYD-1232-SS-257}{比聞蘄州進士馮傑，本儒家，都大坑冶司抑爲爐戶，誅求日增，傑妻以憂死，其女繼之，弟大聲因赴訴，死于道路，傑知不免，毒其二子一妾，舉火自經而死}。

\subsection{1233年：南宋}
《宋史》为这一年留下了可回查的年序入口。

\noindent\textbf{南宋。}\eventanchor{SYD-1233-SS-107}{己-\{丑\}-，詔崔與之、李、鄭性之赴闕}；\eventanchor{SYD-1233-SS-258}{詔袁韶奪職、罷祠祿}。

\subsection{1234年：南宋}
《宋史》为这一年留下了可回查的年序入口。

\noindent\textbf{南宋。}\eventanchor{SYD-1234-SS-108}{甲戌，朱揚祖、林拓朝謁八陵回，以圖進，上問諸陵相去幾何及陵前澗水新復，揚祖悉以對，上忍涕太息}；\eventanchor{SYD-1234-SS-259}{建陽縣盜發，衆數千人，焚劫邵武、麻沙、長平}。

\subsection{1235年：南宋}
《宋史》为这一年留下了可回查的年序入口。

\noindent\textbf{南宋。}\eventanchor{SYD-1235-SS-109}{丙辰，詔主管侍衛馬軍孟珙黃州駐紮，措置邊防}；\eventanchor{SYD-1235-SS-260}{賜進士吳叔告以下四百五十四人及第、出身有差}。

\subsection{1236年：南宋}
《宋史》为这一年留下了可回查的年序入口。

\noindent\textbf{南宋。}\eventanchor{SYD-1236-SS-110}{甲午，詔：「沿江制置使陳韡應援淮東，授淮西制置使兼沿江制置副使史嵩之應援江陵、峽州江面上流}。

\subsection{1237年：南宋}
《宋史》为这一年留下了可回查的年序入口。

\noindent\textbf{南宋。}\eventanchor{SYD-1237-SS-111}{詔曾伯等十一人各官一轉}。

\subsection{1238年：南宋}
《宋史》为这一年留下了可回查的年序入口。

\noindent\textbf{南宋。}\eventanchor{SYD-1238-SS-112}{戊戌，詔：「近覽李奏，知蜀漸次收復，然創殘之餘，綏撫爲急，宜施蕩宥之澤}。

\subsection{1239年：南宋}
《宋史》为这一年留下了可回查的年序入口。

\noindent\textbf{南宋。}\eventanchor{SYD-1239-SS-113}{秉義郎李良守鄂州長壽縣，沒于戰陣，詔贈官三轉}。

\subsection{1240年：南宋}
《宋史》为这一年留下了可回查的年序入口。

\noindent\textbf{南宋。}\eventanchor{SYD-1240-SS-114}{又言：「張祥有保全趙彥呐、楊恢兩制置之功，敵人憚其果毅，宜見錄用}。

\subsection{1241年：南宋}
《宋史》为这一年留下了可回查的年序入口。

\noindent\textbf{南宋。}\eventanchor{SYD-1241-SS-115}{尋以王安石謂「天命不足畏，祖宗不足法，人言不足恤」，爲萬世罪人，豈宜從祀孔子廟庭，黜之}。

\subsection{1242年：南宋}
《宋史》为这一年留下了可回查的年序入口。

\noindent\textbf{南宋。}\eventanchor{SYD-1242-SS-116}{八月丁卯，詔淮東先鋒馬軍鄧淳、李海等揚州撻扒店之戰，宣勞居多，各官兩轉，餘推恩有差}。

\subsection{1243年：南宋}
《宋史》为这一年留下了可回查的年序入口。

\noindent\textbf{南宋。}\eventanchor{SYD-1243-SS-117}{丙辰，安豐軍統領陳友直以王家堈戰功，與官兩轉，壬申，布衣王與之進所著《周禮訂議》，補下州文學}。

\subsection{1244年：南宋}
《宋史》为这一年留下了可回查的年序入口。

\noindent\textbf{南宋。}\eventanchor{SYD-1244-SS-118}{詔<u>玠</u>趣上立功將士姓名等第，即與推恩}。

\subsection{1245年：南宋}
《宋史》为这一年留下了可回查的年序入口。

\noindent\textbf{南宋。}\eventanchor{SYD-1245-SS-119}{詔王雲贈三秩，仍官其二子爲承信郎}。

\subsection{1246年：南宋}
《宋史》为这一年留下了可回查的年序入口。

\noindent\textbf{南宋。}\eventanchor{SYD-1246-SS-120}{詔倪政贈官三轉，官一子承信郎}。

\subsection{1247年：南宋}
《宋史》为这一年留下了可回查的年序入口。

\noindent\textbf{南宋。}\eventanchor{SYD-1247-SS-121}{丙申，以旱，避殿減膳}。

\subsection{1248年：南宋}
《宋史》为这一年留下了可回查的年序入口。

\noindent\textbf{南宋。}\eventanchor{SYD-1248-SS-122}{詔命詞、給告身付之}。

\subsection{1249年：南宋}
《宋史》为这一年留下了可回查的年序入口。

\noindent\textbf{南宋。}\eventanchor{SYD-1249-SS-123}{癸亥，詔給官田五百畝，命臨安府創慈幼局，收養道路遺棄初生嬰兒，仍置藥局療貧民疾病}。

\subsection{1250年：南宋}
《宋史》为这一年留下了可回查的年序入口。

\noindent\textbf{南宋。}\eventanchor{SYD-1250-SS-124}{參知政事謝方叔、吳潛、簽書樞密院事徐清叟並乞解機政，詔不允}。

\subsection{1251年：南宋}
《宋史》为这一年留下了可回查的年序入口。

\noindent\textbf{南宋。}\eventanchor{SYD-1251-SS-125}{丁未，命呂文福廬州駐紮御前諸軍都統制}。

\subsection{1252年：南宋}
《宋史》为这一年留下了可回查的年序入口。

\noindent\textbf{南宋。}\eventanchor{SYD-1252-SS-126}{八月己未，詔來年省試仍舊用二月一日，殿試用四月十五日以前，庶免滯留遠方士子}。

\subsection{1253年：南宋}
《宋史》为这一年留下了可回查的年序入口。

\noindent\textbf{南宋。}\eventanchor{SYD-1253-SS-127}{寶祐元年春正月庚寅朔，詔以藝祖嫡系十一世孫嗣榮王與芮之子建安郡王孜爲皇子，改賜名禥，授崇慶軍節度使，進封永嘉郡王}。

\subsection{1254年：南宋}
《宋史》为这一年留下了可回查的年序入口。

\noindent\textbf{南宋。}\eventanchor{SYD-1254-SS-128}{詔蒲擇之暫權四川制置司事}。

\subsection{1255年：南宋}
《宋史》为这一年留下了可回查的年序入口。

\noindent\textbf{南宋。}\eventanchor{SYD-1255-SS-129}{丙辰，謝方叔、徐清叟以-\{御\}-史朱應元言罷}。

\subsection{1256年：南宋}
《宋史》为这一年留下了可回查的年序入口。

\noindent\textbf{南宋。}\eventanchor{SYD-1256-SS-130}{乙巳，以監察-\{御\}-史吳衍、翁應弼劾太學武學生劉黻等八人不率，詔拘管江西、湖南州軍，宗學生與亻匈等七人並削籍，拘管外宗正司}。

\subsection{1257年：南宋}
《宋史》为这一年留下了可回查的年序入口。

\noindent\textbf{南宋。}\eventanchor{SYD-1257-SS-131}{丙午，禁奸民作白衣會，監司、郡縣官等失覺察者坐罪}。

\subsection{1258年：南宋}
《宋史》为这一年留下了可回查的年序入口。

\noindent\textbf{南宋。}\eventanchor{SYD-1258-SS-132}{丙辰，給事中張鎮言：徐敏子曩帥廣右，嗜殺黷貨，流毒桂府}。

\subsection{1259年：南宋}
《宋史》为这一年留下了可回查的年序入口。

\noindent\textbf{南宋。}\eventanchor{SYD-1259-SS-133}{庚辰，以趙與爲觀文殿學士、兩淮安撫制置使兼知揚州}。

\subsection{1260年：南宋}
《宋史》为这一年留下了可回查的年序入口。

\noindent\textbf{南宋。}\eventanchor{SYD-1260-SS-134}{賈似道奏：「和出彼謀，豈容一切輕徇？倘以交鄰國之道來，當令入見}。

\subsection{1261年：南宋}
《宋史》为这一年留下了可回查的年序入口。

\noindent\textbf{南宋。}\eventanchor{SYD-1261-SS-135}{’請宣付史館，永爲世程法}。

\subsection{1263年：南宋}
《宋史》为这一年留下了可回查的年序入口。

\noindent\textbf{南宋。}\eventanchor{SYD-1263-SS-136}{八月甲寅，董宋臣以病乞收回恩命，請祠，詔賜告五月}。

\subsection{1264年：南宋}
《宋史》为这一年留下了可回查的年序入口。

\noindent\textbf{南宋。}\eventanchor{SYD-1264-SS-137}{丁-\{丑\}-，詔避殿減膳，應中外臣僚許直言朝政闕失}。

\subsection{1265年：南宋}
《宋史》为这一年留下了可回查的年序入口。

\noindent\textbf{南宋。}\eventanchor{SYD-1265-SS-138}{八月庚辰，命陳奕沿江按閱軍防，賜錢二十萬給用}。

\subsection{1266年：南宋}
《宋史》为这一年留下了可回查的年序入口。

\noindent\textbf{南宋。}\eventanchor{SYD-1266-SS-139}{詔自咸淳三年爲始罷之}。

\subsection{1267年：南宋}
《宋史》为这一年留下了可回查的年序入口。

\noindent\textbf{南宋。}\eventanchor{SYD-1267-SS-140}{八月辛酉，遣步帥陳奕率馬軍舟師巡邏江防}。

\subsection{1268年：南宋}
《宋史》为这一年留下了可回查的年序入口。

\noindent\textbf{南宋。}\eventanchor{SYD-1268-SS-141}{於是盧鉞等相繼論列方叔昨蜀、廣敗事，誤國殄民，今又違制擅進，削一秩罰輕}。

\subsection{1269年：南宋}
《宋史》为这一年留下了可回查的年序入口。

\noindent\textbf{南宋。}\eventanchor{SYD-1269-SS-142}{八月戊寅，詔郡縣收民田租，毋巧計取贏，毋厚直折納，轉運司申嚴按劾}。

\subsection{1270年：南宋}
《宋史》为这一年留下了可回查的年序入口。

\noindent\textbf{南宋。}\eventanchor{SYD-1270-SS-143}{癸亥，詔：「贛、吉、南安境數被寇，雖有砦卒，寇出沒無時，莫能相救}。

\subsection{1271年：南宋}
《宋史》为这一年留下了可回查的年序入口。

\noindent\textbf{南宋。}\eventanchor{SYD-1271-SS-144}{詔饒立削兩秩、武岡軍居住}。

\subsection{1272年：南宋}
《宋史》为这一年留下了可回查的年序入口。

\noindent\textbf{南宋。}\eventanchor{SYD-1272-SS-145}{又詔：「有虞之世，三載考績，三考黜陟幽明}。

\subsection{1273年：南宋}
《宋史》为这一年留下了可回查的年序入口。

\noindent\textbf{南宋。}\eventanchor{SYD-1273-SS-146}{帝覽奏，亟詔淮東制司往清口，擇利城築以備之}。

\subsection{1274年：南宋}
《宋史》为这一年留下了可回查的年序入口。

\noindent\textbf{南宋。}\eventanchor{SYD-1274-SS-147}{丙申，江東沙圩租米，以咸淳九年水災，詔減什四}。

\subsection{1275年：南宋}
《宋史》为这一年留下了可回查的年序入口。

\noindent\textbf{南宋。}\eventanchor{SYD-1275-SS-148}{八月己亥朔，總制毛獻忠將衢州兵入衛}。

\subsection{1276年：南宋}
《宋史》为这一年留下了可回查的年序入口。

\noindent\textbf{南宋。}\eventanchor{SYD-1276-SS-149}{八月，漳州亂，命陳文龍爲閩廣宣撫使以討之}。

\subsection{1277年：南宋}
《宋史》为这一年留下了可回查的年序入口。

\noindent\textbf{南宋。}\eventanchor{SYD-1277-SS-150}{八月，文天祥諸將兵皆敗，乃引兵即鄒洬於永豐，洬兵亦潰}。

\subsection{1278年：南宋}
《宋史》为这一年留下了可回查的年序入口。

\noindent\textbf{南宋。}\eventanchor{SYD-1278-SS-151}{曾淵子自雷州來，以爲參知政事，廣西宣諭使}。

\subsection{1279年：南宋}
《宋史》为这一年留下了可回查的年序入口。

\noindent\textbf{南宋。}\eventanchor{SYD-1279-SS-152}{大軍至中軍，會暮，且風雨，昏霧四塞，咫尺不相辨}。
